% Compact drop-in replacement preserving the structure of the original section.
% Numerical note: see real_data_section_revised.tex before final submission.

\section{Real Data Applications}
\label{sec:real-data}

We evaluate the proposed methods on three real regression datasets representing
distinct distributional regimes: heterogeneous wages, a continuous
materials-science outcome, and semicontinuous medical expenditure. For each
dataset, we use 20 repeated random splits, allocating approximately one third of
the observations to training, calibration, and held-out evaluation. The mean,
scale, and conditional-quantile functions are estimated using gradient boosting
with common tuning parameters. We set \(C=0.90\) and \(1-\alpha=0.95\).

Since the conditional response distributions are unknown, we report held-out
empirical content and interval width over repeated random splits. These
quantities are empirical diagnostics and should not be interpreted as direct
estimates of the theoretical conditional PAC guarantees. For CPS and MEPS, the
analysis targets the unweighted public-use sample distribution rather than a
survey-weighted population distribution.

\subsection{CPS Hourly Wage}

The 2012 Current Population Survey sample contains \(29{,}217\) observations,
with hourly wage as the response. The predictors describe sex, marital status,
education, region, and labor-market experience. We include all two-way
interactions among 15 base covariates, yielding 100 nonconstant predictors. The
methods are fitted to log hourly wage, and interval endpoints are transformed
back to dollars per hour.

All four methods attain empirical content above \(C=0.90\) in every split.
SR--TI, ASR--TI, and CQR--TI have average contents \(0.915\), \(0.915\), and
\(0.914\), respectively, compared with \(0.936\) for the Parametric TI.
CQR--TI has the smallest mean width, \(36.32\), while the Parametric TI has the
smallest 90th-percentile width, \(59.66\).

\subsection{Superconducting Critical Temperature}

The superconductivity dataset contains \(21{,}263\) materials represented by 81
physicochemical features. The response is critical temperature in Kelvin,
analyzed on its original scale. The three PAC-calibrated methods have average
empirical content \(0.917\) and attain the target in every split. The Parametric
TI has average content \(0.907\), although its held-out content falls below the
target in 4 of the 20 splits. The Parametric TI has the smallest average width,
while ASR--TI is the narrowest of the PAC-calibrated methods.

\subsection{MEPS Medical Expenditure}

The third application uses Panel 21 of the 2016 Medical Expenditure Panel Survey.
After preprocessing, the analysis contains \(15{,}656\) individuals and 138
predictors. The response is annual total medical expenditure. Expenditure is
zero for \(18.2\%\) of the observations and is strongly right-skewed, so the
methods are fitted to \(\log(1+Y)\) and transformed back to dollars.

The Parametric TI undercovers, with average content \(0.884\). SR--TI, ASR--TI,
and CQR--TI attain average contents \(0.920\), \(0.921\), and \(0.920\),
respectively, and reach the target in every split.
ASR--TI reduces the mean width by approximately \(30.5\%\) relative to SR--TI,
while CQR--TI is the shortest method. Because the response has a point mass at
zero, this application is a stress test of the finite-sample marginal PAC layer
rather than an illustration of the continuous location--scale conditions used
for the conditional SR--TI result.

\begin{table}[H]
\centering
\caption{Real-data empirical content and interval width over 20 repeated random
splits.}
\label{tab:real-data-summary}
\small
\begin{tabular}{llccc}
\toprule
Dataset & Method & Content & Mean width & Q90 width \\
\midrule
CPS & SR--TI          & 0.915 & 37.18 & 63.17 \\
CPS & ASR--TI         & 0.915 & 37.68 & 64.13 \\
CPS & CQR--TI         & 0.914 & \textbf{36.32} & 60.06 \\
CPS & Parametric TI   & 0.936 & 39.77 & \textbf{59.66} \\
\midrule
Supercond. & SR--TI        & 0.917 & 42.61 & 78.61 \\
Supercond. & ASR--TI       & 0.917 & 42.57 & 78.58 \\
Supercond. & CQR--TI       & 0.917 & 45.98 & 79.98 \\
Supercond. & Parametric TI & 0.907 & \textbf{39.41} & \textbf{46.58} \\
\midrule
MEPS & SR--TI        & 0.920 & 80{,}804 & 179{,}317 \\
MEPS & ASR--TI       & 0.921 & 56{,}142 & 125{,}895 \\
MEPS & CQR--TI       & 0.920 & \textbf{13{,}664} & \textbf{29{,}986} \\
MEPS & Parametric TI & 0.884 & 113{,}044 & 313{,}210 \\
\bottomrule
\end{tabular}

\vspace{0.25em}
\begin{flushleft}
\footnotesize
Widths are reported on the original response scale: dollars per hour for CPS,
Kelvin for superconductivity, and dollars for MEPS. Bold entries indicate the
smallest width among methods whose average empirical content exceeds the target
level \(C=0.90\), within each dataset.
\end{flushleft}
\end{table}

\paragraph{Takeaway.}
The three applications illustrate that the proposed methods are validity tools
rather than uniformly shortest intervals. The PAC-calibrated methods attain the
target in every split across all three datasets, whereas the parametric
benchmark is unstable for superconductivity and fails sharply for MEPS.
Efficiency remains score- and data-dependent: CQR--TI is most efficient for CPS
and MEPS, while the Parametric TI is shortest on superconductivity.
